\documentclass[11pt]{article}
\usepackage[letterpaper,margin=0.85in]{geometry}
\usepackage{amsmath,amssymb,booktabs}
\newcommand{\ang}[1]{\langle #1\rangle}
\newcommand{\sq}[1]{[#1]}
\newcommand{\M}{\mathcal M}
\newcommand{\J}{\mathcal J}
\setlength{\parindent}{0pt}
\setlength{\parskip}{0.42em}
\begin{document}
{\LARGE\bfseries Four Fermion Pairs from One Photon Sewing}\par
\textit{Massive tree-level QED in HEPCAT's all-incoming spinor conventions}

Label the electrons $1,3,5,7$ and the positrons $2,4,6,8$. Keep all eight spin indices open.
Let $s_S=(\sum_{i\in S}p_i)^2$. For every unordered division into two nonempty
charge-neutral sets, sew one angle-square spinor tensor on either side:
\begin{equation}
 \boxed{\displaystyle
 \M_8=8e^6\sum_{\{S,\bar S\}}\epsilon(S,\bar S)\,
       \frac{\J(S):\J(\bar S)}{s_S}.}
 \label{eq:sew}
\end{equation}
There are 16 choices of one electron and one positron against the remaining six legs,
and $\binom42^2/2=18$ divisions into two sets of four.
On the $2+6$ terms, evaluate $s_S$ on the two-leg side.
This is particularly useful numerically near a small photon invariant.
The formula serves both pole preferences.

\textbf{The neutral tensor.}
A charged block is a singleton or a triple $ijk$ in which $i,j$ have the same
charge, $i<j$, and $k$ has the opposite charge. For either kind write $W_A$
for its pair of angle and square spinors. Put
\[
 U\odot V=U^\angle\otimes V^\square+V^\angle\otimes U^\square,
 \qquad
 (a^\angle\otimes b^\square):(c^\angle\otimes d^\square)
       =\ang{ca}\sq{bd}.
\]
Extend the contraction by linearity. Everything here carries only spinor indices.
For each neutral $S$, enumerate its divisions into a net-electron block $A$ and a
net-positron block $B$. Then
\[
 \J(S)=\sum_{A\sqcup B=S}\epsilon_S(A,B)\,W_A\odot W_B .
\]
The sign is the parity of the concatenated leg list $(A,B)$ relative to the
increasing list of $S$. Keep $ijk$ in the displayed like-charge-first order.
In Eq.~\eqref{eq:sew}, use the same rule on the two increasing lists $S,\bar S$.
Interchanging the two even-sized lists leaves the sign invariant.
For instance, $\J(12)=W_1\odot W_2$, while
\[
 \J(1234)=-W_1\odot W_{243}+W_3\odot W_{241}
          -W_{132}\odot W_4+W_{134}\odot W_2 .
\]

\textbf{The charged spinor.}
For a singleton, $W_i=(|i\rangle,|i])$. For a triple put $Q=p_i+p_j+p_k$ and
$\tau_{ijk}=s_{ijk}-m^2$. The same rule applies to either net charge:
\begin{align}
 X_i^{jk}&=|j]\ang{ki}+|k]\ang{ji}, &
 Y_i^{jk}&=|k\rangle[ji]+|j\rangle[ki],\nonumber\\
 D_i^{jk}&=\big(QX_i^{jk}-mY_i^{jk},\; QY_i^{jk}-mX_i^{jk}\big),&
 W_{ijk}&=\frac{1}{\tau_{ijk}}
          \left(\frac{D_i^{jk}}{s_{jk}}-\frac{D_j^{ik}}{s_{ik}}\right).
 \label{eq:wave}
\end{align}
$Q$ denotes the corresponding Weyl momentum insertion on each chirality.
The difference in Eq.~\eqref{eq:wave} performs the identical-fermion exchange
inside the triple. Its two physical poles are explicit in each contribution.

\newpage
{\Large\bfseries Why the 34 Sewings Are Complete}

Every eight-fermion tree has six cubic vertices, three internal photons and
two internal electron propagators. Contract each charged line to a vertex.
The resulting tree on four vertices is a star or a path. In the star, the
two ends of the central electron line each attach to a terminal pair.
In the path, each of the two middle lines attaches to one terminal pair.
Cutting both electron propagators thus detaches two disjoint three-leg
branches in either topology. Four charged endpoints remain, two of which
carry those branches. The photon between the two remaining cubic vertices is
unique. Grouping by this photon's external partition gives Eq.~\eqref{eq:sew}.

A terminal photon Fierz contraction supplies the two chiral terms $X$ and
$Y$. Acting with the adjacent electron propagator numerator gives $D$. Combining
its two attachments gives the signed wave $W_{ijk}$. The Fierz contraction at
the middle photon is the colon in Eq.~\eqref{eq:sew}. Each of the three photon
contractions supplies a factor two, giving $8e^6$ in these conventions.
All internal propagators stay off shell.

\begin{center}
\begin{tabular}{@{}lrrr@{}}
\toprule
Middle photon & Channels & Block divisions per channel & Diagrams \\
\midrule
Two legs against six  & 16 & $1\times9$ & $16\cdot9\cdot4=576$ \\
Four legs against four & 18 & $4\times4$ & $18\cdot16\cdot4=1,152$ \\
\bottomrule
\end{tabular}
\end{center}
The last factor four is the product of the two local, two-term fermion exchanges.
The 34 photon channels and 48 possible charged triples also recover the
82 physical channels of the eight-point amplitude. These counts measure
different things from the number of monomials after fully expanding the result.

\textbf{Numerical checks.} At 10,000 phase-space points with $m=0.0005$,
$e=1/3$ and coherent external spins, the maximum relative difference from
the signed sum of 1,728 HEPCAT Feynman diagrams was $3.86\times10^{-12}$
(seed 19300001). At 10,000 different momenta with $m=1$ and the same charge,
the full sum of the 256 squared spin amplitudes agreed with MadGraph to
$1.97\times10^{-13}$ (seed 19600001). A further 1,000 signed checks at
seed 19500001 gave $2.44\times10^{-13}$.

\textbf{Evaluation cost.} Each charged wave or neutral tensor can be tabulated
once for the spin states of its own legs. The final sum retains interference
between all 34 channels at each spin assignment. On an Apple M3, the median
of three 512-point sweeps at seed 19700001 was 0.116 ms per full spin sum
with the standalone sewing evaluator and 3.84 ms with the MadGraph standalone
library. Including phase-space generation and external spinor preparation on
the sewing side raised its time to 0.118 ms. Both timings use the same momenta
and exclude library initialization. The evaluator implements this eight-fermion
formula directly. The general SPINAS exporter was measured separately.

\end{document}
